% Vietnamese translation for OI Public Library
% Source-File: IOI中国国家候选队论文/2006/李天翼/李天翼.doc
\documentclass[11pt]{article}
\usepackage{oipl-vn}

\title{Xét Từ Các Trường Hợp Đặc Biệt}
\author{Li Tianyi, Trường Trung học trực thuộc Đại học Phục Đán}
\date{2006}

\begin{document}
\maketitle

\origtitle{Xét từ các trường hợp đặc biệt}
\sourcepath{IOI China candidate team papers / 2006 / Li Tianyi.doc}

\renewcommand{\abstractname}{Tóm tắt}
\begin{abstract}
Xét từ các trường hợp đặc biệt là một tư tưởng toán học quan trọng. Các trường
hợp đặc biệt chủ yếu gồm hai loại: trường hợp đơn giản và trường hợp cực đoan.

Bài viết dùng một vài ví dụ để minh họa cách áp dụng tư tưởng này trong thi
tin học, đồng thời rút ra điểm chung của chúng và chỉ ra nội hàm quan trọng
của tư tưởng đó.
\end{abstract}

\noindent\textbf{Từ khóa:} trường hợp đặc biệt, thi tin học.

\tableofcontents

\section{Ví dụ 1: Bra}

\subsection{Mô tả bài toán}

Bài này được phóng tác từ POI 2003--2004, bài \textit{Bra}.

Xét một mạch gồm \(n\) cổng. Các cổng được đánh số
\(0,1,2,\ldots,n-1\). Mỗi cổng có một số lượng đầu vào cố định và một đầu ra.
Đầu vào và đầu ra có thể nhận một trong ba trạng thái \(0,1,\frac12\). Mỗi đầu
vào nối với đầu ra của một cổng nào đó. Trạng thái của đầu vào bằng trạng thái
của đầu ra mà nó nối tới. Mỗi đầu ra có thể nối tới nhiều đầu vào. Hai cổng
được đánh số \(0\) và \(1\) là đặc biệt: chúng không có đầu vào; cổng \(0\)
luôn xuất ra \(0\), còn cổng \(1\) luôn xuất ra \(1\).

Ta nói trạng thái đầu ra của một cổng là \term{hợp lệ} khi và chỉ khi thỏa một
trong các điều kiện sau:
\begin{enumerate}
  \item nó bằng \(0\) và trong các đầu vào của cổng đó, số \(0\) nhiều hơn số
  \(1\);
  \item nó bằng \(\frac12\) và trong các đầu vào của cổng đó, số \(0\) bằng số
  \(1\);
  \item nó bằng \(1\) và trong các đầu vào của cổng đó, số \(1\) nhiều hơn số
  \(0\);
  \item nó bằng số hiệu của cổng, và số hiệu đó là \(0\) hoặc \(1\).
\end{enumerate}

Nếu trạng thái đầu ra của mọi cổng đều hợp lệ, ta gọi mạch là \term{hợp lệ}.
Nếu trạng thái đầu ra của một cổng là như nhau trong mọi mạch hợp lệ, ta nói
trạng thái đầu ra của cổng đó là cố định. Dữ liệu bảo đảm tồn tại ít nhất một
mạch hợp lệ.

\paragraph{Nhiệm vụ.}
Hãy viết một chương trình thực hiện các việc sau:
\begin{itemize}
  \item đọc mô tả mạch từ đầu vào chuẩn;
  \item với mỗi cổng, kiểm tra trạng thái đầu ra của nó có cố định hay không;
  nếu có, xác định trạng thái đó;
  \item ghi trạng thái của các cổng có đầu ra cố định ra đầu ra chuẩn.
\end{itemize}

\paragraph{Dữ liệu vào.}
Đầu vào chuẩn chứa một số nguyên \(n\), \(2\le n\le 10000\). Tiếp theo là
\(n-2\) dòng mô tả kết nối của từng cổng. Dòng thứ \(i\) mô tả các đầu vào của
cổng \(i\): số nguyên đầu tiên là \(k_i\), \(k_i\ge 1\), cho biết cổng này có
\(k_i\) đầu vào; tiếp theo là \(k_i\) số hiệu cổng tương ứng với các đầu ra
được nối vào. Các số nguyên trên cùng một dòng cách nhau bằng dấu cách. Tổng số
đầu vào của mọi cổng không vượt quá \(200000\).

\paragraph{Dữ liệu ra.}
Chương trình phải ghi \(n\) dòng ra đầu ra chuẩn. Dòng thứ \(i\) phụ thuộc vào
trạng thái đầu ra của cổng được đánh số \(i-1\):
\begin{itemize}
  \item \texttt{0}: nếu nó luôn bằng \(0\);
  \item \texttt{1/2}: nếu nó luôn bằng \(\frac12\);
  \item \texttt{1}: nếu nó luôn bằng \(1\);
  \item \texttt{?}: nếu nó không xác định.
\end{itemize}

\paragraph{Ví dụ.}
\begin{verbatim}
Dữ liệu vào:
5
2 0 1
2 4 2
2 2 4

Dữ liệu ra:
0
1
1/2
?
?
\end{verbatim}

\subsection{Lời giải}

Do đồ thị có chu trình, ta khó trực tiếp phán đoán trạng thái đầu ra của từng
cổng có cố định hay không. Đây chính là trở ngại của bài toán.

Gọi \(P(i)\) là trạng thái đầu ra của cổng \(i\), với
\(0\le i\le n-1\). Gọi \(P_{\min}(i)\) và \(P_{\max}(i)\) lần lượt là giá trị
nhỏ nhất và lớn nhất mà \(P(i)\) có thể nhận trong mọi mạch hợp lệ; đó là các
trường hợp cực đoan của \(P(i)\).

Rõ ràng, nếu
\[
  P_{\min}(i)=P_{\max}(i),\qquad 0\le i\le n-1,
\]
thì trạng thái đầu ra của cổng \(i\) là cố định; ngược lại thì không cố định.
Vì vậy ta chỉ cần tìm \(P_{\min}(i)\) và \(P_{\max}(i)\).

Ký hiệu \(C_{j,i}\) là số đầu vào của cổng \(i\) nối với đầu ra của cổng \(j\).
Xét giá trị
\[
  \frac{\sum_{j=0}^{n-1} C_{j,i}P(j)}
       {\sum_{j=0}^{n-1} C_{j,i}},
\]
tức trung bình các trạng thái đầu vào của cổng \(i\), với
\(2\le i\le n-1\). Theo định nghĩa ``hợp lệ'' trong đề bài, trong mọi mạch hợp
lệ:
\begin{itemize}
  \item nếu giá trị này nhỏ hơn \(\frac12\), thì \(P(i)=0\);
  \item nếu giá trị này bằng \(\frac12\), thì \(P(i)=\frac12\);
  \item nếu giá trị này lớn hơn \(\frac12\), thì \(P(i)=1\).
\end{itemize}

Ta thực hiện thao tác sau. Ban đầu đánh dấu trạng thái đầu ra của mọi cổng là
\(0\); khi đó chỉ cổng \(1\) chưa hợp lệ. Bắt đầu từ cổng \(1\), đổi trạng thái
đầu ra của nó thành \(1\). Sau đó liên tục tìm nơi phát sinh mâu thuẫn và lặp
việc cập nhật.

Ta chứng minh quá trình lặp này chắc chắn dừng, và khi dừng thì
\[
  P(i)=P_{\min}(i),\qquad 0\le i\le n-1.
\]

Giả sử mệnh đề không đúng. Ở thời điểm bắt đầu thao tác, với mọi
\(i\), \(0\le i\le n-1\), ta có \(P(i)\le P_{\min}(i)\). Vì mệnh đề không
đúng, chắc chắn có một thời điểm nào đó bắt đầu xuất hiện
\(P(k)>P_{\min}(k)\). Ngay trước thời điểm đó, với mọi
\(i\), \(0\le i\le n-1\), vẫn có \(P(i)\le P_{\min}(i)\).

Xét trung bình các trạng thái đầu vào của cổng \(k\):
\[
  \frac{\sum_{j=0}^{n-1} C_{j,k}P(j)}
       {\sum_{j=0}^{n-1} C_{j,k}}.
\]
Giá trị này đã đủ lớn để việc lấy \(P(k)=P_{\min}(k)\) không còn thỏa yêu cầu.
Mặt khác,
\[
  \frac{\sum_{j=0}^{n-1} C_{j,k}P(j)}
       {\sum_{j=0}^{n-1} C_{j,k}}
  \le
  \frac{\sum_{j=0}^{n-1} C_{j,k}P_{\min}(j)}
       {\sum_{j=0}^{n-1} C_{j,k}}.
\]
Điều này có nghĩa là không tồn tại một mạch hợp lệ nào thỏa
\(P(k)=P_{\min}(k)\), mâu thuẫn với định nghĩa của \(P_{\min}(k)\). Do đó mệnh
đề được chứng minh.

Mỗi cổng nhiều nhất đổi trạng thái hai lần, từ \(0\) sang \(\frac12\) và từ
\(\frac12\) sang \(1\). Tổng số đầu vào của mọi cổng không vượt quá
\(200000\), nên sau không quá \(2\cdot 200000=400000\) lần lặp, quá trình sẽ
dừng. Khi đó \(P(i)=P_{\min}(i)\), với \(0\le i\le n-1\).

Tương tự, ta có thể tìm \(P_{\max}(i)\), \(0\le i\le n-1\). Đến đây toàn bộ
bài toán được giải quyết.

\subsection{Tiểu kết}

Trường hợp cực đoan là một dạng biểu hiện của trường hợp đặc biệt. Nhiều tính
chất trong đề bài thường bộc lộ qua các đối tượng có tính cực đoan, chẳng hạn
việc lấy cực trị trong bài này. Chính vì vậy ta có thể tập trung khảo sát
chúng để tìm điểm đột phá và đáp án.

\section{Ví dụ 2: Sko}

\subsection{Đặt vấn đề}

\subsubsection{Mô tả bài toán}

\textit{Sko} là bài của POI 2004--2005.

Một quân mã di chuyển trên bàn cờ vô hạn. Mỗi loại nước đi mà nó có thể thực
hiện được biểu diễn bởi một cặp số nguyên. Cặp \((a,b)\) nghĩa là quân mã có
thể đi từ điểm có tọa độ \((x,y)\) tới \((x+a,y+b)\) hoặc tới
\((x-a,y-b)\). Mỗi quân mã có một tập gồm nhiều cặp số nguyên; các cặp đó biểu
diễn toàn bộ các nước đi mà quân mã này có thể thực hiện. Với mỗi quân mã, giả
sử các điểm nó đi tới được từ gốc \((0,0)\) không nằm tất cả trên cùng một
đường thẳng.

Ta nói hai quân mã là \term{giống nhau} nếu tập các điểm chúng có thể đi tới
từ \((0,0)\), sau số bước tùy ý và không nhất thiết bằng nhau giữa hai quân,
là hoàn toàn giống nhau. Có thể biết rằng với mỗi quân mã, luôn có một quân mã
``giống'' nó và chỉ cần hai cặp số nguyên để mô tả.

\paragraph{Nhiệm vụ.}
Hãy viết chương trình:
\begin{itemize}
  \item đọc các cặp số nguyên biểu diễn nước đi của quân mã;
  \item xác định hai cặp số nguyên mô tả một quân mã ``giống'' quân mã đã cho;
  \item ghi hai cặp số nguyên đó ra đầu ra chuẩn.
\end{itemize}

\paragraph{Dữ liệu vào.}
Dòng đầu của đầu vào chuẩn chứa một số nguyên \(n\), là số cặp số nguyên
\((3\le n\le 100)\). Trong \(n\) dòng tiếp theo, mỗi dòng chứa một cặp số
nguyên biểu diễn một loại nước đi của quân mã. Trên các dòng này, hai số
\(a_i\) và \(b_i\) cách nhau bởi một dấu cách. Ta có
\(-100\le a_i,b_i\le 100\), và giả sử \((a_i,b_i)\ne(0,0)\).

\paragraph{Dữ liệu ra.}
Dòng đầu của đầu ra chuẩn chứa hai số nguyên \(a,b\) cách nhau bởi một dấu
cách. Dòng thứ hai chứa hai số nguyên \(c,d\) cách nhau bởi một dấu cách.
Các số nguyên này thỏa \(-10000\le a,b,c,d\le 10000\), và quân mã có hai nước
đi \((a,b)\), \((c,d)\) phải ``giống'' quân mã được mô tả trong dữ liệu vào.

\paragraph{Ví dụ.}
\begin{verbatim}
Dữ liệu vào:
3
24 28
15 50
12 21

Dữ liệu ra:
468 1561
2805 9356

hoặc:
3 0
0 1
\end{verbatim}

\subsubsection{Ý tưởng ban đầu}

Muốn xét quân mã đã cho ``giống'' với loại quân mã nào, trước hết cần biết nó
có thể đi tới những điểm nào. Ta gọi các điểm mà một quân mã có thể đi tới từ
\((0,0)\) là các \term{điểm khả thi} của quân mã đó. Tập các điểm khả thi của
một quân mã gọi là \term{tập điểm khả thi}.

Bàn cờ là hai chiều. Ta có thể bắt đầu từ trường hợp một chiều đơn giản hơn.
Khi đó mỗi quân mã chỉ di chuyển trên một đường thẳng; nếu nó có \(n\) loại
nước đi, nước đi thứ \(i\), \(1\le i\le n\), có thể biểu diễn bằng một số
nguyên \(a_i\). Đặt
\[
  r=\gcd(a_1,a_2,\ldots,a_n).
\]
Điều kiện cần và đủ để quân mã đi tới được điểm có hoành độ \(X\) là
\[
  r\mid X.
\]

Với trường hợp hai chiều không suy biến, tức các điểm có thể đi tới không nằm
tất cả trên một đường thẳng, có thể phỏng đoán tập điểm khả thi của nó trùng
với tập mọi điểm thỏa một trong ba dạng điều kiện sau:
\[
  r\mid X+Y,\qquad
  r\mid X+qY,\qquad
  r\mid pX+qY,
\]
trong đó \(p,q,r\) đều là các số nguyên cần xác định.

Xét trường hợp \(n=2\), \((a_1,b_1)=(2,3)\), \((a_2,b_2)=(6,6)\). Dễ thấy hai
giả thiết đầu là sai. Với giả thiết cuối, vì có nhiều tham số hơn nên trong
nhất thời khó phán đoán nó đúng hay sai.

\subsection{Hai thuật toán chuẩn bị}

\subsubsection{Thuật toán Euclid}

Ký hiệu ước chung lớn nhất của hai số nguyên không âm \(a,b\) là
\(\gcd(a,b)\). Phương pháp thông dụng nhất để tìm ước chung lớn nhất là thuật
toán Euclid, dựa trên định lý sau.

\paragraph{Định lý đệ quy GCD.}
Với mọi số nguyên không âm \(a\) và mọi số nguyên dương \(b\),
\[
  \gcd(a,b)=\gcd(b,a\bmod b).
\]

Trước hết chứng minh \(\gcd(a,b)\mid \gcd(b,a\bmod b)\). Đặt
\(d=\gcd(a,b)\), khi đó \(d\mid a\) và \(d\mid b\). Ta có
\[
  a\bmod b = a-qb,\qquad q=\left\lfloor\frac ab\right\rfloor.
\]
Vì \(d\mid a\) và \(d\mid qb\), nên \(d\mid a-qb\), tức
\(d\mid a\bmod b\). Do đó \(\gcd(a,b)\mid \gcd(b,a\bmod b)\).

Ngược lại, chứng minh \(\gcd(b,a\bmod b)\mid \gcd(a,b)\). Đặt
\(d=\gcd(b,a\bmod b)\), khi đó \(d\mid b\) và \(d\mid a-qb\), suy ra
\(d\mid a\). Vì vậy \(\gcd(b,a\bmod b)\mid \gcd(a,b)\). Hai chiều chia hết cho
nhau, nên \(\gcd(a,b)=\gcd(b,a\bmod b)\).

Mã giả của thuật toán Euclid như sau:
\begin{verbatim}
Euclid(a,b)
  if b=0
     then return a
     else return Euclid(b,a mod b)
\end{verbatim}

Xét số lần gọi đệ quy của thủ tục Euclid. Giả sử \(a>b\ge 0\); nếu
\(b>a\ge 0\), sau một bước ta sẽ gọi \(\textsc{Euclid}(b,a)\), còn nếu
\(b=a>0\), thủ tục chỉ thực hiện một lần.

Đặt
\[
  F_k=\frac1{\sqrt5}\left[
    \left(\frac{1+\sqrt5}{2}\right)^k
    -
    \left(\frac{1-\sqrt5}{2}\right)^k
  \right],
\]
với \(k\) là số nguyên dương bất kỳ.\footnote{Trong bản PDF,
dấu giữa hai lũy thừa trong công thức Binet bị hỏng; ở đây dùng dạng chuẩn của
số Fibonacci, phù hợp với hệ thức \(F_{k+3}=F_{k+2}+F_{k+1}\) được dùng ngay
sau đó.}
Ta chứng minh: nếu \(a>b\ge 0\) và \(\textsc{Euclid}(a,b)\) thực hiện
\(n\ge 1\) lần gọi đệ quy, thì
\[
  a\ge F_{n+2},\qquad b\ge F_{n+1}.
\]

Với \(n=1\), ta có \(b\ge 1=F_2\), và \(a>b\) nên \(a\ge 2=F_3\); mệnh đề hiển
nhiên đúng. Giả sử mệnh đề đúng với \(n=k\), trong đó \(k\ge 1\). Xét
\(n=k+1\). Vì \(k>0\), nên \(b>0\), và \(\textsc{Euclid}(a,b)\) gọi đệ quy.
Theo giả thiết quy nạp,
\[
  b\ge F_{k+2},\qquad a\bmod b\ge F_{k+1}.
\]
Ta có
\[
  b+(a\bmod b)
  = b+\left(a-\left\lfloor\frac ab\right\rfloor b\right)
  \le a.
\]
Vì vậy
\[
  a\ge b+(a\bmod b)\ge F_{k+2}+F_{k+1}=F_{k+3}.
\]
Do đó mệnh đề cũng đúng với \(n=k+1\).

Từ mệnh đề trên suy ra: với mọi số nguyên \(k\ge 1\), nếu \(a>b\ge 0\) và
\(b<F_{k+1}\), thì số lần gọi đệ quy của \(\textsc{Euclid}(a,b)\) nhỏ hơn
\(k\). Vì \(F_k\) tăng theo cấp số mũ, độ phức tạp thời gian của thuật toán
Euclid bị chặn bởi \(O(\lg(\min(a,b)))\).

\subsubsection{Cách giải phương trình tuyến tính modulo}

\paragraph{Định nghĩa.}
Phương trình tuyến tính modulo là phương trình có dạng
\[
  ax\equiv n\pmod b,
\]
trong đó \(a,b,x,n\) đều là số nguyên và \(ab\ne 0\).

\paragraph{Nghiệm.}
Xét phương trình tuyến tính modulo tổng quát
\[
  ax\equiv n\pmod b.
\]
Nó có thể chuyển thành phương trình
\[
  ax+by=n.
\]
Nếu \(n\) không chia hết cho \(\gcd(a,b)\), phương trình hiển nhiên vô nghiệm.
Nếu \(n\) chia hết cho \(\gcd(a,b)\), đặt \(d=\gcd(a,b)\) và xét phương trình
đặc biệt
\[
  ax\equiv d\pmod b.
\]
Khi đó phương trình có thể viết thành \(ax+by=d\). Đặt
\[
  a'=\frac ad,\qquad b'=\frac bd.
\]
Theo định lý Bézout, tồn tại hai số nguyên \(p,q\) sao cho
\[
  a'p+b'q=1,
\]
do đó phương trình này có nghiệm. Đặt
\[
  x_0=\frac{pn}{d}.
\]
Khi đó \(ax_0\equiv n\pmod b\), tức \(x_0\) là một nghiệm riêng của phương
trình.

Đặt
\[
  x=x_0+\frac{mb}{d},
\]
với \(m\) là số nguyên bất kỳ, thì vẫn có \(ax\equiv n\pmod b\). Ngoài ra,
với mọi \(x_1\), nếu \(ax_1\equiv n\pmod b\), thì
\[
  b\mid a(x_1-x_0),
\]
suy ra
\[
  \frac bd\mid x_1-x_0.
\]

\paragraph{Tìm một nghiệm riêng.}
Một nghiệm riêng của \(ax\equiv d\pmod b\) có thể được tìm bằng mã giả sau:
\begin{verbatim}
Extended_Euclid(a,b)
  if b=0
     then return (a,1,0)
  (d',p',q') = Extended_Euclid(b,a mod b)
  (d,p,q) = (d',q',p' - floor(a/b) * q')
  return (d,p,q)
\end{verbatim}

Với phương trình tuyến tính modulo tổng quát \(ax\equiv n\pmod b\), mã giả sau
cho một nghiệm riêng:
\begin{verbatim}
Modular_Linear_Equation_Solver(a,n,b)
  (d,x',y') = Extended_Euclid(a,b)
  if d | n
     then x = x' * (n/d) mod b
     else no solution
\end{verbatim}

Dễ thấy thuật toán này cũng có độ phức tạp thời gian
\(O(\lg(\min(a,b)))\).

\subsection{Giải bài toán}

\subsubsection{Chứng minh phỏng đoán}

\paragraph{Đặc biệt hóa phỏng đoán.}
Ký hiệu \(S(p,q,r)\) là tập mọi điểm thỏa
\[
  r\mid pX+qY,
\]
trong đó \(p,q,r\) đều là số nguyên và \(r\ne 0\).

Xét phỏng đoán ban đầu, cũng là mệnh đề cần chứng minh: với trường hợp hai
chiều không suy biến, tức các điểm quân mã có thể đi tới không nằm tất cả trên
một đường thẳng, tồn tại các số nguyên \(p,q,r\), \(r\ne 0\), sao cho tập điểm
khả thi của quân mã trùng với \(S(p,q,r)\).

Khi \(n=1\), bài toán suy biến thành một chiều, nên không thỏa phỏng đoán. Vì
giới hạn trên của \(n\) là \(100\), khá lớn và không tiện thảo luận trực tiếp,
ta xét trước trường hợp đơn giản hơn \(n=2\), cũng là trường hợp không suy
biến nhỏ nhất. Phần dưới cho thấy chỉ cần xử lý \(n=2\).

Giả sử phỏng đoán đúng khi \(n=k\), với \(k\) là số nguyên dương bất kỳ. Khi
\(n=k+1\), giả sử quân mã chỉ dùng \(k\) nước đi đầu có tập điểm khả thi trùng
với \(S(p,q,r_0)\). Với quân mã có đủ \(k+1\) nước đi, \((X,Y)\) là điểm khả
thi khi và chỉ khi tồn tại số nguyên \(m\) sao cho
\[
  p(X-ma_{k+1})+q(Y-mb_{k+1})\equiv 0\pmod{r_0}.
\]
Điều này tương đương với
\[
  (pa_{k+1}+qb_{k+1})m\equiv pX+qY\pmod{r_0}.
\]
Xét phương trình tuyến tính modulo này, ta biết nó có nghiệm khi và chỉ khi
\[
  \gcd(pa_{k+1}+qb_{k+1},r_0)\mid pX+qY.
\]
Đặt \(r=\gcd(pa_{k+1}+qb_{k+1},r_0)\), suy ra phỏng đoán cũng đúng với
\(n=k+1\). Vì vậy nếu giải được \(n=2\), thì mọi trường hợp \(n\ge 2\) đều
được giải.

Khi \(n=2\), quân mã chỉ có hai nước đi \((a_1,b_1)\) và \((a_2,b_2)\), với
\[
  a_1b_2-a_2b_1\ne 0,
\]
nếu không bài toán suy biến thành một chiều.

Để đào sâu bản chất bài toán rõ hơn, trước hết ta chứng minh một mệnh đề đơn
giản hơn.

\paragraph{Mệnh đề 1.}
Quân mã có hai nước đi \((a_1,b_1)\) và \((a_2,b_2)\), thỏa
\[
  a_1b_2-a_2b_1\ne 0,\qquad \gcd(a_1,a_2)=1,\qquad \gcd(b_1,b_2)=1.
\]
Khi đó tồn tại các số nguyên \(p,q,r\), \(r\ne 0\), sao cho tập điểm khả thi
của quân mã trùng với \(S(p,q,r)\).

Khác biệt bản chất giữa Mệnh đề 1 và mệnh đề gốc là Mệnh đề 1 bổ sung điều
kiện \(\gcd(a_1,a_2)=1\) và \(\gcd(b_1,b_2)=1\).

\paragraph{Chứng minh xây dựng của Mệnh đề 1.}
Theo định lý Bézout, tồn tại các số nguyên \(x,y\) sao cho
\[
  a_1x+a_2y=1.
\]
Chọn hai số nguyên \(s_1,s_2\) thỏa
\[
  a_1s_1+a_2s_2=1.
\]
Đặt
\[
  p=b_1s_1+b_2s_2,\qquad q=-1,\qquad r=a_1b_2-a_2b_1.
\]

Trước hết chứng minh tính cần thiết: nếu \((X,Y)\) là điểm khả thi của quân
mã, thì \(r\mid pX+qY\). Ta có
\[
\begin{aligned}
  pa_1+qb_1
  &= (b_1s_1+b_2s_2)a_1-b_1 \\
  &= (a_1s_1-1)b_1+b_2s_2a_1 \\
  &= -a_2s_2b_1+a_1s_2b_2
   = s_2(a_1b_2-a_2b_1)=s_2r.
\end{aligned}
\]
Tương tự,
\[
  pa_2+qb_2=-s_1r.
\]
Vì \((X,Y)\) là điểm khả thi, giả sử quân mã thực hiện nước đi
\((a_1,b_1)\) tổng cộng \(c\) lần và nước đi \((a_2,b_2)\) tổng cộng \(d\) lần
rồi đi từ gốc tới \((X,Y)\). Khi đó
\[
\begin{aligned}
  pX+qY
  &=p(ca_1+da_2)+q(cb_1+db_2)\\
  &=c(pa_1+qb_1)+d(pa_2+qb_2)\\
  &=cs_2r-ds_1r
   =r(cs_2-ds_1).
\end{aligned}
\]
Do đó \(r\mid pX+qY\).

Tiếp theo chứng minh tính đủ: nếu \(r\mid pX+qY\), thì \((X,Y)\) là điểm khả
thi của quân mã. Vì \(r\mid pX+qY\), tồn tại số nguyên \(k\) sao cho
\[
  pX+qY=kr.
\]
Do \(q=-1\), ta có \(pX-Y=kr\). Mặt khác,
\[
\begin{aligned}
  X
  &=X\cdot 1+0\\
  &=X(a_1s_1+a_2s_2)+(ka_1a_2-ka_1a_2)\\
  &=(Xs_1+ka_2)a_1+(Xs_2-ka_1)a_2,
\end{aligned}
\]
và
\[
\begin{aligned}
  Y
  &=pX+(Y-pX)=Xp-kr\\
  &=X(b_1s_1+b_2s_2)-k(a_1b_2-a_2b_1)\\
  &=(Xs_1+ka_2)b_1+(Xs_2-ka_1)b_2.
\end{aligned}
\]
Vì vậy nếu quân mã thực hiện nước đi \((a_1,b_1)\) tổng cộng
\(Xs_1+ka_2\) lần và nước đi \((a_2,b_2)\) tổng cộng \(Xs_2-ka_1\) lần, nó sẽ
đi từ gốc tới \((X,Y)\). Mệnh đề 1 được chứng minh.

\paragraph{Trở lại phỏng đoán.}
Bây giờ, để chứng minh phỏng đoán ban đầu, chỉ cần chuyển nó về Mệnh đề 1.
Điểm đặc biệt của Mệnh đề 1 là \(a_1,a_2\) nguyên tố cùng nhau và
\(b_1,b_2\) nguyên tố cùng nhau.

Đặt
\[
  d_1=\gcd(a_1,a_2),\qquad d_2=\gcd(b_1,b_2).
\]
Xét quân mã có hai nước đi
\[
  \left(\frac{a_1}{d_1},\frac{b_1}{d_2}\right),\qquad
  \left(\frac{a_2}{d_1},\frac{b_2}{d_2}\right).
\]
Theo Mệnh đề 1, tồn tại các số nguyên \(p,q,r\), \(r\ne 0\), sao cho tập điểm
khả thi của quân mã này trùng với \(S(p,q,r)\).

Đặt \(d=\gcd(p,q)\). Chỉ cần xét trường hợp \(d=1\); nếu không, thay
\[
  p\leftarrow \frac pd,\qquad q\leftarrow \frac qd,\qquad r\leftarrow \frac rd.
\]
Ta phải có \(\gcd(p,r)=1\), nếu không thì tung độ của mọi điểm trong
\(S(p,q,r)\) đều là bội của \(\gcd(p,r)\), mâu thuẫn với điều kiện
\(\gcd(b_1/d_2,b_2/d_2)=1\).

Đặt
\[
  m=\gcd(r,d_1),\qquad d_1=mt.
\]
Phương trình bất định
\[
  rx\equiv 1-p\pmod t
\]
chắc chắn có nghiệm nguyên. Gọi \(x_0\) là một nghiệm của phương trình này,
và đặt
\[
  p'=rx_0+p.
\]
Khi đó \(\gcd(p',m)=1\), \(\gcd(p',t)=1\), nên \(\gcd(p',d_1)=1\).

Tương tự, ta có thể thu được \(q'\). Với \(p'\) và \(q'\) thỏa
\[
  \gcd(p',d_1)=1,\qquad \gcd(q',d_2)=1,
\]
có thể biết rằng quân mã có hai nước đi \((a_1,b_1)\), \((a_2,b_2)\) có tập
điểm khả thi là
\[
  S(p'd_2,\ q'd_1,\ rd_1d_2).
\]
Đến đây phỏng đoán ban đầu đã được chứng minh.

\subsubsection{Từ chứng minh tới thuật toán}

\paragraph{Một thủ tục phụ.}
Từ chứng minh phỏng đoán, ta có thể biểu diễn tập điểm khả thi của một quân mã
chỉ có hai nước đi bằng \(S(p,q,r)\). Sau đó, dựa trên chứng minh xây dựng cho
trường hợp \(n>2\), phạm vi biểu diễn được mở rộng tới quân mã bất kỳ. Tuy
nhiên, bài toán yêu cầu xuất một quân mã chỉ có hai nước đi nhưng có tập điểm
khả thi là \(S(p,q,r)\). Ta cần các nước đi của quân mã đó.

Đặt \(d=\gcd(p,q)\). Có thể giả sử \(d=1\); nếu không, thay
\[
  p\leftarrow \frac pd,\qquad
  q\leftarrow \frac qd,\qquad
  r\leftarrow \frac rd.
\]
Đặt
\[
  d_1=\gcd(p,r),\qquad
  d_2=\gcd(q,r),\qquad
  p'=\frac p{d_1},\qquad
  q'=\frac q{d_2},\qquad
  r'=\frac r{d_1d_2}.
\]
Gọi \(t\) là một nghiệm của phương trình tuyến tính modulo
\[
  q't\equiv -1\pmod{r'}.
\]
Khi đó hai nước đi
\[
  (0,d_2r')\qquad\text{và}\qquad (d_1,d_2t)
\]
thỏa yêu cầu. Tính đúng đắn có thể rút ra từ chứng minh của Mệnh đề 1.

\paragraph{Thủ tục chính.}
Tư tưởng cơ bản của thuật toán là chuyển tập điểm khả thi của quân mã sang
dạng \(S(p,q,r)\), rồi chuyển dạng này thành tập điểm khả thi của một quân mã
chỉ có hai nước đi.

Gọi \(P\) là giá trị tuyệt đối lớn nhất trong mọi \(a_i,b_i\)
\((1\le i\le n)\). Độ phức tạp thời gian của thuật toán là
\[
  O(n\lg |P|).
\]

\subsection{Tiểu kết}

Trong một số bài toán, liên hệ giữa điều kiện và kết quả không rõ ràng, khiến
ta khó tìm điểm đột phá. Khi xuất phát từ nơi nguyên thủy nhưng vẫn quan trọng
nhất, chẳng hạn trường hợp một chiều trong bài này, ta có thể nhìn rõ vấn đề.
Hoặc ta dùng một trường hợp đơn giản, như Mệnh đề 1, làm tấm gương để so sánh
với trường hợp tổng quát, tức phỏng đoán trong bài này; từ đối chiếu đó tìm ra
khác biệt bản chất nhất giữa hai tình huống, rồi xử lý đúng chỗ để giải quyết
triệt để bài toán.

Đặc biệt trong một số bài toán dựng hình hoặc xây dựng nghiệm, có rất nhiều
trường hợp khả thi nhưng ta chỉ cần một trong số đó. Khi ấy chỉ cần xét một
trường hợp tương đối đơn giản, như bước chuyển từ Mệnh đề 1 sang mệnh đề gốc
trong bài này, là có thể giải quyết nhanh và hiệu quả. Những ví dụ như vậy có
rất nhiều, chẳng hạn Ví dụ 3.

\section{Ví dụ 3: Polygon}

\subsection{Mô tả bài toán}

\textit{Polygon} là bài của BOI 2005.

Hãy tìm một đa giác lồi sao cho mọi cạnh của nó có các độ dài đã cho.

\paragraph{Dữ liệu vào.}
Tên tệp là \texttt{poly.in}. Dòng đầu của tệp chứa một số nguyên
\(n\), \(3\le n\le 1000\), là số cạnh của đa giác lồi. \(n\) dòng tiếp theo,
mỗi dòng chứa một số nguyên \(a_i\), \(1\le a_i\le 10000\), biểu diễn độ dài
của một cạnh.

\paragraph{Dữ liệu ra.}
Nếu đa giác lồi như vậy tồn tại, tệp \texttt{poly.out} phải chứa \(n\) dòng.
Mỗi dòng chứa hai số thực \(x_i,y_i\), với
\(|x_i|,|y_i|\le 10000000\). Nối \((x_i,y_i)\) với
\((x_{i+1},y_{i+1})\) cho \(1\le i\le n-1\), và nối \((x_n,y_n)\) với
\((x_1,y_1)\); đa giác thu được phải có các độ dài cạnh đúng bằng các độ dài
đã cho, nhưng không bắt buộc theo cùng thứ tự.

Thứ tự các điểm xuất ra có thể theo chiều kim đồng hồ hoặc ngược chiều kim
đồng hồ. Nếu đa giác lồi như vậy không tồn tại, chỉ cần xuất
\texttt{NO SOLUTION}.

\paragraph{Ví dụ.}
\begin{verbatim}
Dữ liệu vào:
4
7
4
5
4

Dữ liệu ra:
0.5 2.5
7.5 2.5
4.5 6.5
0.5 6.5
\end{verbatim}

\paragraph{Chấm điểm.}
Nếu trị tuyệt đối của hiệu giữa hai số thực không vượt quá \(0.001\), trình
chấm coi chúng là bằng nhau. Mọi định dạng xuất số thực đều được chấp nhận.

\subsection{Giải bài toán}

Vì giữa các cạnh không có thứ tự, ta có thể giả sử
\[
  a_1\ge a_2\ge a_3\ge\cdots\ge a_n.
\]
Rõ ràng, nếu
\[
  \sum_{i=2}^{n} a_i\le a_1,
\]
thì vô nghiệm; ngược lại thì có nghiệm. Phần dưới chỉ thảo luận trường hợp có
nghiệm.

\subsubsection{Một ý tưởng ngây thơ}

Chọn tùy ý một điểm ban đầu \((x_1,y_1)\) rồi nối các cạnh ra ngoài. Khi nối
cạnh thứ \(k\), giả sử \(d\) là khoảng cách giữa \((x_1,y_1)\) và
\((x_{k+1},y_{k+1})\). Ta yêu cầu
\[
  d<\sum_{i=k+1}^{n} a_i
  \qquad\text{và}\qquad
  a_{k+1}<d+\sum_{i=k+2}^{n} a_i.
\]
Khi đó cuối cùng chắc chắn có thể quay về điểm xuất phát.

Tuy nhiên, cách xử lý này chỉ bảo đảm mỗi đoạn của đường gấp khúc khép kín có
độ dài đúng như đã cho; nó không bảo đảm đường đó không tự cắt, càng không bảo
đảm nó là đa giác lồi.

\subsubsection{Xét trường hợp đặc biệt}

Ta xét một trường hợp tương đối đặc biệt: nếu đa giác lồi \(n\) cạnh này là
một đa giác nội tiếp đường tròn thì sao?

Rõ ràng mọi đa giác nội tiếp đường tròn đều lồi. Vì vậy nếu có thể dựng một
đa giác \(n\) cạnh nội tiếp đường tròn với mỗi cạnh có độ dài đúng như đã cho,
đa giác đó chắc chắn thỏa yêu cầu. Do đó chỉ cần tìm đường kính của đường tròn
này.

Giả sử đường kính \(d\) đã biết. Lấy \((0,0)\) làm tâm đường tròn, lấy điểm
\((d/2,0)\) làm điểm xuất phát, rồi dựng đa giác lồi \(n\) cạnh theo chiều kim
đồng hồ.

Khi cạnh \(a_1\) và \(a_n\) chưa giao nhau thì \(d\) quá lớn; khi \(a_1\) và
\(a_n\) cắt chéo thì \(d\) quá nhỏ. Vì miền giá trị của \(d\) là liên tục,
đường kính \(d\) thỏa điều kiện chắc chắn tồn tại và có thể tìm bằng tìm kiếm
nhị phân.

Giả sử đường kính đường tròn ngoại tiếp của đa giác là \(d\). Khi đó góc ở tâm
tương ứng với cạnh thứ \(i\) là
\[
  2\arcsin\frac{a_i}{d}.
\]
Hiển nhiên \(d\ge a_1\).

Xét quan hệ giữa
\[
  \sum_{i=2}^{n}2\arcsin\frac{a_i}{d}
\]
và \(\pi\). Có ba trường hợp:
\begin{enumerate}
  \item Hai giá trị bằng nhau. Khi đó đặt \(d=a_1\) là đủ thỏa yêu cầu.
  \item Tổng phía trước lớn hơn. Khi đó đường tròn thỏa yêu cầu có tâm nằm
  trong đa giác. Đường kính \(d\) thỏa
  \[
    \sum_{i=1}^{n}2\arcsin\frac{a_i}{d}=2\pi.
  \]
  Vì góc ở tâm ứng với cạnh dài nhất không nhỏ hơn \(2\pi/n\), nên
  \[
    d\le \frac{a_1}{\sin(\pi/n)}.
  \]
  \item Tổng phía sau lớn hơn.\footnote{Bản PDF ghi ``tổng phía sau
  lớn hơn'', nhưng theo ngữ cảnh đây là trường hợp tổng
  \(\sum_{i=2}^{n}2\arcsin(a_i/d)\) nhỏ hơn \(\pi\). Tôi giữ ý nghĩa được suy
  ra từ lập luận hình học.} Khi đó đường tròn thỏa yêu cầu có tâm nằm ngoài đa
  giác. Đường kính \(d\) thỏa
  \[
    \sum_{i=2}^{n}2\arcsin\frac{a_i}{d}
    =2\arcsin\frac{a_1}{d}.
  \]
  Lúc này \(d\) có thể rất lớn. Tuy nhiên, do giới hạn dữ liệu
  \[
    n\le 1000,\qquad a_i\le 10000,\qquad a_i\text{ đều là số nguyên},
  \]
  trên thực tế giá trị lớn nhất của \(d\) chỉ khoảng \(5\cdot 10^5\).
\end{enumerate}

Đến đây bài toán đã được giải quyết.

\section{Tổng kết}

Đối diện với thế giới rộng lớn, ta thường chỉ nhìn thấy biểu tượng tổng thể
bên ngoài của sự vật, rồi dựa vào cảm giác của mình mà đi ngày càng xa. Ta
không biết rằng chính mình bị biểu tượng đó che mắt, nên không thể tìm ra bí
mật ẩn bên trong.

Pythagoras cho rằng ``vạn vật đều là số''. Xét từ các trường hợp đặc biệt thực
chất là một tư tưởng toán học; nó có thể giúp ta, trong nhiều phương diện kể
cả thi tin học, nhận thức bản chất của sự vật tốt hơn.

Khi đã nhận ra bản chất của sự vật, ta có thể giải quyết vấn đề tương đối dễ
dàng. Hơn nữa, trong quá trình giải quyết vấn đề, ta còn thu được gợi ý hữu
ích để cung cấp đối sách cho các tình huống tương tự. Điều này quan trọng hơn
bản thân việc giải xong một bài toán.

\section*{Tài liệu tham khảo}
\addcontentsline{toc}{section}{Tài liệu tham khảo}

\begin{enumerate}
  \item Thomas H. Cormen, Charles E. Leiserson, Ronald L. Rivest, Clifford
  Stein, \textit{Introduction to Algorithms}, 2nd edition.
  \item Wu Wenhu và Wang Jiande, \textit{Phân tích thuật toán thực dụng và lập
  trình}.
  \item Su Chun, \textit{Nhìn vấn đề từ tính đặc thù}.
\end{enumerate}

\end{document}
